\section{Validation: Bi\texorpdfstring{$_2$}{2}Se\texorpdfstring{$_3$}{3} and Bi\texorpdfstring{$_2$}{2}Te\texorpdfstring{$_3$}{3} films}
\label{sec:results_ordered}

Ordered Bi$_2$Se$_3$ and Bi$_2$Te$_3$ films provide validation. Detector images collected at incident angles of $5^\circ$, $10^\circ$, and $15^\circ$ entered the staged refinement. Figures~\ref{fig:ordered_bi2se3_test} and~\ref{fig:ordered_bi2te3_test} show the $5^\circ$ condition because it displays both the selected off-specular rod families and the $r=0$ features clearly within one detector image while the others can be found in the supplemental materials.

$r\neq0$ rods are naturally described by $Q$ space and integrated along $Q_R$ as a function of $L$. Mosaic and beam sampling distort the projection to $\mathbf Q$ and thus the plotted curves therefore compare the same detector-derived observable rather than an ideal reciprocal-space line cut. Meanwhile, $r=0$ rods are described best integrated along $\phi$ as a function of $2\theta$.

In the detector panels, solid colored curves mark calculated trajectory centerlines. Semi-transparent bands show the regions integrated for the profiles, and dashed outlines mark their boundaries. The central blue region follows the $r=0$ family. Paired off-specular regions are the two detector-side intersections of one $r$ family, identified by plus and minus signs. In the profile panels, black solid curves are measured data and orange dashed curves are the corresponding calculations. The horizontal coordinate is $L$, with integer markers denoting nominal Bragg conditions along the selected trajectory. Off-specular rows are plotted on a linear scale, while the $r=0$ row is plotted logarithmically to display weak off-condition intensity. Along the $r=0$ trajectory, the radial $L$ or $2\theta$ line shape contains the wavelength convolution, whereas the transverse $\phi$ width provides the cleaner mosaic constraint.

\begin{figure}[p]
  \centering
  \captionsetup{font=footnotesize}
  {\large\bfseries Measured detector image and selected trajectory bands\par}
  \vspace{0.35em}
  \begin{minipage}{0.78\linewidth}
    \centering
    \draftfigureplaceholder[0.28\textheight]{figures/results_ordered/figure7_bi2se3_qr_rod_qz_profiles_detector_selected_q_regions_5deg.png}{%
      Bi$_2$Se$_3$ detector image with selected fixed-$Q_R$ trajectory centerlines and finite profile-integration regions.
    }
  \end{minipage}

  \vspace{0.35em}
  {\large\bfseries Matched projected profiles\par}
  \vspace{0.25em}
  \begin{minipage}{0.56\linewidth}
    \centering
    \draftfigureplaceholder[0.35\textheight]{figures/results_ordered/figure7_bi2se3_qr_rod_qz_profiles.png}{%
      Bi$_2$Se$_3$ measured and calculated profiles for the selected detector regions.
    }
  \end{minipage}
  \caption[Bi2Se3 detector-space test case]{Bi$_2$Se$_3$ detector-space comparison at $\theta_i=5^\circ$. The upper panel is the measured detector image with calculated trajectory centerlines, finite integration regions, and their boundaries. Labels $r$ identify radial rod families, while plus and minus signs distinguish the two detector-side intersections of one off-specular family. The lower panel compares measured black curves with calculated orange dashed curves generated using the same coordinate map, integration limits, and area normalization. The figure tests trajectory placement and apparent profile line shape. In the $r=0$ family, radial broadening is evaluated with the wavelength distribution, while the transverse $\phi$ width is used to constrain mosaicity. Separate family amplitudes are not a common-scale relative-intensity test, and the $r=0$ row does not establish a unique broad-component functional form.}
  \label{fig:ordered_bi2se3_test}
\end{figure}

\begin{figure}[p]
  \centering
  \captionsetup{font=footnotesize}
  {\large\bfseries Measured detector image and selected trajectory bands\par}
  \vspace{0.35em}
  \begin{minipage}{0.78\linewidth}
    \centering
    \draftfigureplaceholder[0.28\textheight]{figures/results_ordered/figure7_bi2te3_qr_rod_qz_profiles_detector_selected_q_regions_5deg.png}{%
      Bi$_2$Te$_3$ detector image with selected fixed-$Q_R$ trajectory centerlines and finite profile-integration regions.
    }
  \end{minipage}

  \vspace{0.35em}
  {\large\bfseries Matched projected profiles\par}
  \vspace{0.25em}
  \begin{minipage}{0.56\linewidth}
    \centering
    \draftfigureplaceholder[0.35\textheight]{figures/results_ordered/figure7_bi2te3_qr_rod_qz_profiles.png}{%
      Bi$_2$Te$_3$ measured and calculated profiles for the selected detector regions.
    }
  \end{minipage}
  \caption[Bi2Te3 detector-space test case]{Bi$_2$Te$_3$ detector-space comparison at $\theta_i=5^\circ$, using the conventions of Fig.~\ref{fig:ordered_bi2se3_test}. The upper panel shows the measured detector image with calculated trajectories and finite integration regions. The lower panel compares detector-derived profiles formed with the same coordinate transformation, finite integration, and area normalization. The principal peaks occur at the same projected locations and have similar apparent line shapes within the displayed fitted condition. Separate family amplitudes are not interpreted on one common physical scale, and the $r=0$ profile is not proof of a unique broad mosaic component.}
  \label{fig:ordered_bi2te3_test}
\end{figure}

Mosaicity, beam divergence, wavelength spread, incidence-angle uncertainty, calibrated detector mapping, and finite integration width all broaden the effective scattering condition. The projected coordinate should therefore not be interpreted as an exact cut through the reciprocal space of one perfect crystallite. The comparison remains direct because the same coordinate transformation and finite integration operation are applied to the calculated and measured detector quantities.

The off-specular and $r=0$ features play complementary roles. The off-specular arcs test the two-dimensional-powder reciprocal-space construction and calibrated detector geometry. Their transverse widths and asymmetries constrain the effective narrow mosaic component together with beam divergence and finite integration width. For the $r=0$ family, the radial $2\theta$ or $L$ shape is sensitive to the wavelength distribution and the specular intensity envelope, while the transverse $\phi$ width is the cleaner probe of the broad orientation component. The full-spectrum versus mean-wavelength comparison in Sec.~\ref{sec:correlated_effects} isolates the radial $003$ streak from the mosaic-sensitive transverse width. The displayed overlays show that the fitted model can represent the selected feature locations and apparent profile shapes. A controlled Gaussian-only comparison is still required before the broad component can be called necessary or uniquely Lorentzian.

